\documentclass{article}
\usepackage{amsmath,amssymb}
\usepackage[margin=1in]{geometry}

\title{Where the solution spaces touch\\
\large CANDIDATE model, not a theorem}
\author{Note on \texttt{j15marti/pnp}, branch \texttt{main}, identity \texttt{j15marti}.\\
Status of the whole model: CANDIDATE. Nothing here is claimed as proved.}
\date{2026-10-04, about 12:20~PM PT}

\begin{document}
\maketitle

\noindent
P versus NP remains open.
This note writes one candidate model of how a polynomial region and a forbidden region sit in one coordinate, which touch can actually be shown, and which limit points are only a formula or a toy.
It does not solve P versus NP, it does not assert \(\mathrm{P}=\mathrm{NP}\) or \(\mathrm{P}\neq\mathrm{NP}\), and it does not exhibit an algorithm that decides the question.
It does not edit root \texttt{main.tex}.
It does not add a workflow.

The principle is taken from the already committed note
\texttt{research/agents/grok/2026-10-04\_1207\_GROK\_SQUEEZE1.tex}
(commit \texttt{69f6a8fdbaa78920a74b493dbb97354af76d2bd6}, blob \texttt{1a61162fecc531487803439983e05b2507b423be}), which was read for this note.
The finite charges in
\texttt{research/agents/grok/2026-10-04\_1206\_GROK\_PC\_RAM2.tex}
(commit \texttt{24375c54c9264cdc9975c5ced679d1010c986477}, blob \texttt{511b4a4ec11e1b0b2c068a2abadded02075937fd})
are used only as a computational illustration.
They are not a limit.
The literature statements below are those the squeeze note recorded from texts it opened.
This note does not re-fetch those papers and does not add sources.

\section{Common coordinate}
Let \(T(n)>1\) and let the time exponent be
\[
e_T(n)=\frac{\log T(n)}{\log n},
\]
any fixed base greater than \(1\), as in the squeeze note.
Compactify by
\[
z=\frac{1}{1+e_T}.
\]
A finite constant exponent \(k\) sits at the positive point \(z=1/(1+k)\).
If \(e_T(n)\to\infty\), then \(z\to 0\).

The common coordinate of the model is this \(z\), together with a space exponent \(s\in[0,\infty]\) when a statement needs space.
A point with both coordinates finite is a finite resource vector.
The symbol \(\infty\) is the end of the compactification, not a measured runtime.

\section{The two sets}
The polynomial set is the set of finite exponents in that coordinate:
\[
P_{\mathrm{res}}
=
\left\{\frac{1}{1+k}:k\in[0,\infty)\right\}
\subset(0,1].
\]
Equivalently, in the exponent plane, \(P_{\mathrm{res}}\) is the set of pairs \((t,s)\) in which the time exponent \(t\) is finite.
There is no largest polynomial: \(\mathrm{P}\) is the union of \(O(n^k)\) over finite \(k\), not one \(O(n^K)\) for a maximum \(K\).
Superpolynomial rates are those with unbounded exponent, hence those that land at \(z=0\).

The forbidden set \(F\) is not that face \(z=0\).
\(F\) is only the region the squeeze note records from opened theorems, and only with the restrictions those theorems state.
It is not claimed to be every known lower bound, and it is not a region this note measured.

\begin{itemize}
\item Williams, as recorded from the opened extended abstract arXiv:1001.0746.
The time lower bound cited there for SAT on \(n^{o(1)}\)-space algorithms is \(n^{2\cos(\pi/7)-o(1)}\).
The same record states that the paper's comparison writes this as at least \(n^{1.801}\).
That decimal is the paper's comparison, as the squeeze note already says.
It is not recomputed here.
The same record states Theorem~3.9: if time and space are both bounded by polynomials, then any SAT algorithm satisfies \(t\cdot s=\Omega(n^{2\cos(\pi/7)-\varepsilon})\) for every \(\varepsilon>0\).
So, inside finite exponents, a polynomial pair \((c,d)\) with \(c+d<2\cos(\pi/7)\) is a piece of \(F\).
Since \(\cos(\pi/7)<1\), one has \(2\cos(\pi/7)<2\), and the finite pair \((1,1)\) has sum \(2\), so it is outside that piece.
\(F\) does not contain every finite exponent.
\item The same record: this exponent is the limit of a recurrence on alternation-trading proof annotations under subpolynomial space, and Buss--Williams (ECCC TR11-031, abstract only in that note) is cited as saying \(c=2\cos(\pi/7)\) is optimal for the \(\varepsilon=0\) case of that proof system.
That is a boundary of a space-restricted region and of a proof calculus.
It is not a map sending SAT algorithms to SAT algorithms.
\item Fortnow--Lipton--van Melkebeek--Viglas, as recorded from the opened author-page abstract of the JACM 2005 paper, not from a re-opened PDF.
For any constant \(c\) less than the golden ratio there is a positive \(d\) such that no deterministic random-access Turing machine solves SAT in time \(n^c\) and space \(n^d\), and \(d\) approaches \(1\) as \(c\) does.
On conondeterministic machines the same shape holds for any constant \(c\) less than \(\sqrt{2}\).
No decimal for the golden ratio or for \(\sqrt{2}\) is given on that page, and none is introduced here.
This is a space-restricted piece of \(F\).
It is not an algorithm map.
\end{itemize}

\section{Touch}
The touch that this coordinate actually shows is with the superpolynomial boundary, not with \(F\).

The set \(P_{\mathrm{res}}\) accumulates at \(0\): every neighborhood of \(z=0\) contains \(1/(1+k)\) for all sufficiently large finite \(k\).
Every superpolynomial example in the squeeze note, \(n^{\log\log n}\) and \(2^{\sqrt{n}}\), also satisfies \(z\to 0\).
Therefore
\[
\overline{P_{\mathrm{res}}}\cap\{z=0\}=\{0\}
\]
in the compactification.
The closure of the polynomial set meets the superpolynomial boundary at the single point \(z=0\).
That point is an accumulation point of finite exponents.
Taking limits in \(z\) does not put a polynomial rate on one side of a gap and a superpolynomial rate on the other.
The squeeze note already records this, and it is a remark about this coordinate, not a separation of \(\mathrm{P}\) from \(\mathrm{NP}\).

What is not shown is
\[
\overline{P_{\mathrm{res}}}\cap F\neq\emptyset
\]
as a touch that would separate the classes.
The opened set \(F\) cuts a space-restricted piece out of the finite exponents: time below \(2\cos(\pi/7)\) only under \(n^{o(1)}\) space, the product corner \(c+d<2\cos(\pi/7)\) only as a constraint on pairs, and the Fortnow--Lipton--van Melkebeek--Viglas region only with a positive space exponent tied to the time exponent.
Large finite exponents stay outside those pieces.
Their closure still contains \(z=0\), and \(z=0\) is not shown to lie in \(F\).
A limit of distinct polynomial space exponents toward a small \(s\) is not, by itself, one algorithm that uses space \(n^{o(1)}\).

\section{Limit points}
Each point below is labeled \textbf{EXAMPLE}, \textbf{KNOWN}, or \textbf{OPEN}.
No unlabeled constant is a SAT result.

\begin{center}
\begin{tabular}{@{}llp{2.7in}@{}}
\hline
Point & Label & What it is \\
\hline
\(z=0\) & KNOWN &
Shown touch of \(\overline{P_{\mathrm{res}}}\) with the superpolynomial boundary.
Not a separation. \\
\(2\cos(\pi/7)\) & KNOWN &
Williams boundary of the \(n^{o(1)}\)-space region, and of the polynomial product constraint.
Not a contraction fixed point. \\
golden ratio & KNOWN &
FLMV deterministic boundary.
Space-restricted: some \(d>0\), with \(d\to 1\) as \(c\) approaches the golden ratio from below.
No decimal. \\
\(\sqrt{2}\) & KNOWN &
FLMV conondeterministic boundary.
Still space-restricted.
No decimal. \\
\(a/(1-b)\) & OPEN &
Attractor formula only.
No \(a\), \(b\) are known for SAT. \\
\((I-M)^{-1}b\) & OPEN &
Vector formula from the squeeze note.
No matrix \(M\) and no offset \(b\) are constructed. \\
\(2\), from \(a=1\), \(b=1/2\) & EXAMPLE &
Toy fixed point in the next section.
Not a SAT exponent. \\
\hline
\end{tabular}
\end{center}

The paper's comparison \(n^{1.801}\) is not a separate limit point.
It is the decimal the squeeze note attributes to the opened Williams text, and it is not recomputed here.

\section{Toy contraction, labeled EXAMPLE}
The following numbers were computed for this note.
They are not taken from a SAT algorithm.

Choose \(a=1\) and \(b=1/2\), and choose the toy starting exponent \(K_0=6\).
The affine step \(K\mapsto a+bK\) is the map \(K\mapsto 1+K/2\).
The closed form
\[
K_t
=
a\frac{1-b^t}{1-b}+b^t K_0
=
2+4\cdot 2^{-t}
\]
and the iteration \(K_{t+1}=1+K_t/2\) agree:
\[
\begin{align*}
K_0&=6,\\
K_1&=4,\\
K_2&=3,\\
K_3&=5/2,\\
K_4&=9/4,\\
K_5&=17/8,\\
K_6&=33/16,\\
K_7&=65/32.
\end{align*}
\]
The limit is
\[
\frac{a}{1-b}=\frac{1}{1-1/2}=2.
\]
\textbf{EXAMPLE.}
This is a contraction on a toy exponent.
It is not a SAT theorem, not a lower bound, and not an upper bound.
No such constants \(a\) and \(b\) are known for SAT.
The map that would send every exact polynomial SAT algorithm of exponent \(K\) to another exact SAT algorithm of exponent at most \(a+bK\), with the same \(a\) and \(b\) at every iteration, every construction cost charged, and fixed point inside \(F\), is the map the squeeze note marks \textbf{OPEN}.
It remains open here.

\section{Finite charges are not a limit}
The following figures are the computational illustration in the PC\_RAM2 note.
They describe one finite ordered diagram and one frozen formula.
They are not an asymptotic boundary, not a point of \(\overline{P_{\mathrm{res}}}\), and not a point of \(F\).

\begin{itemize}
\item The sign-round-robin ordered clause-query DAG for \(n=4\), \(m=8\), promise weight \(8\), has \(720\) nonterminal nodes and peak level width \(66\).
\item A \(7\)-bit level-local node index, with that DAG as the program, matched the mask-union answer on all \(\binom{32}{8}=10518300\) weight-\(8\) inputs.
\item The direct mutable charge of \(20\) bits (a \(16\)-bit mask plus a \(4\)-bit selected-count, program counter excluded) is slack against the live width: \(20\neq\log_2 66\) and \(20\neq 7\).
\item On one frozen weight-\(8\) CNF, the monomial-space model in that note has minimum space \(5\), and the recomputed resolution width is \(3\).
The inequality \(5>3\) is for that formula in those two models.
It is not a limit of exponents.
\end{itemize}

\section{SAME\_REALIZATION}
A limit of one encoding and a limit of another encoding are not the same object, and are not a contradiction, unless a theorem says they are the same realization.

\section{What would still have to be proved}
A separation by contraction would need all of the following, and none of them is proved here.
A uniform, cost-charged, repeatable map on exact SAT algorithms, with spectral contraction in the scalar form \(0<b<1\) or in the vector form \(\rho(M)<1\).
A fixed point \(a/(1-b)\), or \((I-M)^{-1}b\), lying in the forbidden set \(F\) of one and the same realization, not merely in a toy copy of the recurrence and not merely in the finite \(720\)-node diagram.
An unconditional theorem that really forbids that fixed point, rather than a space-restricted boundary that large finite exponents never enter.
Williams' exponent \(2\cos(\pi/7)\) does not supply that map.
The Fortnow--Lipton--van Melkebeek--Viglas region does not supply that map.

\section{Conclusion}
P versus NP remains open.
In the compactification \(z=1/(1+e_T)\), the closure of the polynomial set meets the superpolynomial boundary at \(z=0\), and that is the touch this model can show.
It cannot show that the same closure meets the forbidden set \(F\) in a way that separates the classes.
The named boundaries \(2\cos(\pi/7)\), the golden ratio, and \(\sqrt{2}\) are boundaries of space-restricted regions.
The attractor \(a/(1-b)\) is a formula.
The only numerical fixed point computed here is the toy value \(2\).
No SAT contraction is known whose fixed point lies in \(F\).
\end{document}
